%! Author = lili
%! Date = 6/9/26
% \chapter{VW3 - Translation}

% BIs F 5 ist Wdh.
\section*{Vorwissen}
\defin{orf}

\begin{figure}
    % TODO F 9
    \caption{Open Reading Frame}
\end{figure}

\section{Funktionsformen der RNA}
% todo F 6 - 7

\defin{mrna}
\subsection{tRNA}
\defin{trna}
% TODO trna F 10 - 17
\defin{kleeblatt}
\begin{figure}
    % TODO F 11
    \caption{Kleeblatt-Struktur}
\end{figure}

\defin{dihydrouridine}
\defin{pseudouridine}

\begin{figure}
    % TODO F 13
    \caption{Pseudouridine und Dihydrouridine vs. Uridin}
\end{figure}

\defin{aminoacyl_trna}
\defin{aminoacyl_trna_synthasen}

\begin{figure}
    % TODO F 14 - 15
    \caption{Aminoacyl-tRNA-Synthetasen}
\end{figure}

% TODO weitere Rechnerhce Codon - tRNA
% TODO Bild F 17

\subsubsection{Funktion der tRNA}
% TODO F 20

\subsubsection{Codon-Anticodon Wechselwirkung}
\defin{codonanticodon}
% TODO F 19

\section{Die Aminosäuren}
 % TODO tabelle F 21
% TODO Karteikarten
% TODO eigenschaften der hydrophoben. basischen etc. Aminosäuren
\frage{wievieleaminoinproteinen}
% TODO F 51-42
% + TODO richtig in karteikarten

\subsection{rRNA}
\defin{rrna}
\begin{figure}
    % TODO F 23
    \caption{Funktionelle Einheiten des Ribosoms}
\end{figure}

% TODO ? \defin{Peptidyl-tRNA} ?? F 23

% TODO Karteikarten Aufbau Ribosom
% TODO Karteikarten Aufgaben der Untereinheiten

\begin{figure}
    % TODO F 25
    \caption{Aufbau des
    prokaryotischen Ribosoms}
\end{figure}

% TODO Karteikarten Aufbau prok Ribosom

\section{Translation}
\defin{translation}
\subsection{Initiation}
% TODO F 29 - 36 
\subsection{Elongation}
% TODO F 37 - 43 
\subsection{Termination}
% TODO F 44 - 47

% TODO was los mit F 48 - 49? 

\subsection{Räumliche Kopplung in Prokaryoten}
% TODO F 50
